\section{未来方向与开放问题}

\subsection{长期依赖与任务组合}

Q-Learning 原始形式在处理长 horizon/ sparse reward 时会变得低效。现代方法倾向于多步 Q 更新（multi-step TD）、hierarchical policy 以及外部 planner 的组合。

\begin{importantbox}{任务组合的 pipeline}
通过把单一任务拆成 Macro / Micro 两层（如 meta policy 负责子任务选择，micro policy 负责具体控制）可以缓解 long-term credit assignment 的困难。
\end{importantbox}

\subsection{治理与可迁移性}

\begin{knowledgebox}{开放问题清单}
\begin{itemize}
    \item Sim-to-real gap：如何让 replay buffer 中真实数据发挥更大作用。
    \item Reward hacking：长期 reward shaping 可能引入 shortcut，需定期复查 reward formula。
    \item 多模态状态：当 robot state 包含图像、LiDAR 与语言，Q network 的处理 pipeline 复杂度骤增。
\end{itemize}
\end{knowledgebox}

\begin{warningbox}{部署时的 drift 管理}
即便策略在 offline 模拟中表现优异，Observation drift 或 action constraint drift 也会导致线上失效。建议配置周期性的 sanity check（如 evaluate on replayed expert rollouts）。
\end{warningbox}

\subsection{本章小结}

本章指出 Q-Learning 领域仍在演进，长期依赖、任务组合与治理审计是未来工作的主要方向。

